% 本文件由 LaTeX 排版 Agent 根据有效 translation.json 自动生成。
% author=Miscellaneous; work_id=0854; title=宗徒的训诲

\chapter{宗徒的训诲}

% structure: p0001 source=h1 target=omitted reason=page-title-already-rendered-by-parent

那时，基督被接升到祂的圣父那里；宗徒们如何领受了圣神的恩赐；教会的规条与法律；每位宗徒去了何方；罗马帝国领土内的地区从何处领受了司铎的祝圣。\par

在希腊王国纪元三百三十九年，\OriginalTerm{赫齐兰月}{Heziran}，同月第四天，即一周的第一天，五旬节的末尾——就在那一天，门徒们从加里肋亚的\Place{纳匝肋}（吾主降孕宣告之处）来到那称为橄榄山的地方，吾主与他们同在，但对他们不可见。在黎明时分，吾主举起双手，按在十一位宗徒的头上，将司铎职的恩赐赐予他们。忽然，一片光明的云彩接走了祂。他们看见祂升天而去。祂坐在圣父的右边。他们赞美天主，因为他们看见祂的升天正如祂曾告诉他们的；他们欢喜，因为他们领受了赋予他们\Person{梅瑟}和\Person{亚郎}家室司铎职的\OriginalTerm{右手}{Right Hand}。\par

从那里他们上到\Emph{城里}，进了一间楼房——就是吾主与他们一起守逾越节的那间，也是曾问询过谁将出卖吾主给钉死者的地方；那里也曾\Emph{有}问询：他们应如何在世上宣讲祂的福音？并且，正如在这楼房内吾主身体与血的奥迹开始在世界得胜，同样，从那里，祂宣讲的教导开始在世上拥有权柄。\par

当门徒们陷入这种困惑，不知如何向言语陌生又不为他们所知的民族宣讲祂的福音，彼此这样说：虽然我们确信基督要借我们的手在言语陌生的民族面前行大能和奇迹，而他们也不通晓我们的语言，然而谁去教导他们，使他们明白这些大能和奇迹是因被钉的基督之名而行呢？——我说，当门徒们正思虑这些事时，\Person{西满}\OriginalTerm{刻法}{Cephas}站起来，对他们说：弟兄们，我们如何宣讲祂的福音，这事不关乎我们，而关乎我们的主；因为\Emph{祂}知道我们如何在世上宣讲祂的福音；我们信赖祂对我们的眷顾，祂曾应许我们说：\Quote{当我升到我父那里去时，我要派遣圣神，那护慰者，\Emph{祂}要教导你们一切应当知道和宣告的事。}\par

当\Person{西满}\OriginalTerm{刻法}{Cephas}向他的同伴宗徒说这些话，提醒他们时，他们听到一种神秘的声音，并有一种世上奇异、甘甜的香气吹向他们；在声音与香气之间，\OriginalTerm{火舌}{tongues of fire}自天而降，落在他们每人身上；每人按各自所领受的方言，便准备前往那讲那方言的地方。\par

并且，借着同一圣神的恩赐（那天赐给他们的），他们制定了\OriginalTerm{规条与法律}{Ordinances and Laws}——合乎他们福音宣讲及真实忠信教导的教义：——\par

1. 宗徒们于是制定：你们要向东祈祷，因为\Quote{闪电从东方发出，直照到西方，人子的来临也要这样。}\BibleRef{玛 24:27}好使我们借此知道并明白祂将从东方突然显现。\par

2. 宗徒们进一步制定：在一周的第一\Emph{日}，要有礼仪、圣经诵读和\OriginalTerm{祭献}{oblation}；因为在一周的第一天，吾主从死者中复活，在一周的第一天祂向世界显现，在一周的第一天祂升天，在一周的第一天祂最后将与众天使一同显现。\par

3. 宗徒们进一步制定：在一周的第四\Emph{天}，要有礼仪；因为在那一天，吾主向他们透露了关于祂的受审、受苦、被钉、死亡和复活的事；门徒们因此忧伤。\par

4. 宗徒们进一步制定：在\Emph{安息日}前夕，第九时辰，要有礼仪；因为在第四天所说关于救主受难的事，在那\Emph{同一}前夕实现了；世界和受造物战栗，天上的\OriginalTerm{光体}{luminaries}变暗。\par

5. 宗徒们进一步制定：要有长老和执事，如同肋未人；以及\OriginalTerm{副执事}{subdeacons}，如同在主的圣所庭院中搬运器皿的人；并有一位\OriginalTerm{监督}{overseer}，他要同样作全民的\OriginalTerm{引导者}{Guide}，如同\Person{亚郎}，整个城中一切司祭和肋未人的首领和元首。\par

6. 宗徒们进一步制定：庆祝我们救主的主显节，这是教会节日的首领，在\OriginalTerm{后卡努月}{latter Canun}第六日，按希腊长历。\par

7. 宗徒们进一步制定：在救主受难日之前四十天，你们要禁食，然后庆祝受难日和复活日；因为我们的主自己，这节日的主，曾禁食四十天；\Person{梅瑟}和\Person{厄里亚}，他们都曾领受这奥秘，也同样各禁食四十天，然后受到光荣。\par

8. 宗徒们进一步制定：在所有\Emph{其他}圣经诵读结束时，要诵读福音，作为一切圣经的\OriginalTerm{印证}{seal}；百姓要站立聆听，因为这是所有人的救赎福音。\par

9. 宗徒们进一步制定：在祂复活后满五十天时，你们要纪念祂升到荣耀的圣父那里。\par

10. 宗徒们制定：除了旧约、先知书、福音和（他们的）宗徒大事录外，在教会内的\OriginalTerm{读经台}{pulpit}上不可读其它。\par

11. 宗徒们进一步规定：凡不熟悉教会信德以及其中所定的规章和法律者，不得担任指导者和治理者；凡熟悉这些却背离它们者，不得再次任职：因为他在任职上不忠实，便说了谎。\par

12. 宗徒们进一步规定：凡发誓、说谎、作假见证、求问巫术、占卜者和加色丁人，相信命运和星象——那些不认识天主的人所执着的——这样的人，也作为不认识天主者，应被逐出任职，不得\Emph{再次}任职。\par

13. 宗徒们进一步规定：若有人\Emph{心意不定}地对待任职，并不以坚定不移的意志追随它，不得再次任职：因为任职之主并未被他以坚定不移的意志事奉；他\Emph{只}欺骗了人，而不是天主，\Quote{在他面前诡计无效。}\par

14. 宗徒们进一步规定：凡放贷收取利息、从事买卖和贪得无厌者，不得再次任职，也不得继续任职。\par

15. 宗徒们进一步规定：凡喜爱犹太人——像喜爱他们的朋友依斯加略那样——或喜爱外邦人（他们崇拜受造物而非造物主）者，不得进入他们中间并任职；而且，若他\Emph{已}在他们中间，他们不应容忍他\Emph{留下}，而应将他从他们中间分离，不得再与他们一同任职。\par

16. 宗徒们进一步规定：若有人从犹太人或外邦人中前来与他们联合，且在联合之后又转回他\Emph{先前}所站的一方，然后又再次返回来到他们那里，则不应再接纳他；而应按他先前的立场，认识他的人怎样看他，就怎样看他。\par

17. 宗徒们进一步规定：指导者不得未经与他一同任职者商议，处理属于教会的事务；而应凭他们所有人的意见发令，\Emph{只}应实行他们全体一致同意且不反对的事。\par

18. 宗徒们进一步规定：凡有人带着对基督信德的美好见证，并因祂的名忍受苦难而离开此世，你们应在他们被处死的那一日纪念他们。\par

19. 宗徒们进一步规定：在教会的礼仪中，你们要日复一日地重复达味的赞美诗：因为\Emph{这话}：\Quote{我要时时赞美主，祂的赞美常在我口中；}以及\Emph{这话}：\Quote{我昼夜默想、讲论，使我的声音在袮面前被听见。}\par

20. 宗徒们进一步规定：凡摆脱玛门、不追逐钱财之利者，应被拣选并接纳到祭台的任职中。\par

21. 宗徒们进一步规定：若司铎偶然违背正义将\Emph{他人}捆绑，应接受应有的惩罚；而被捆绑者应视那捆绑如同公正的捆绑。\par

22. 宗徒们进一步规定：若惯常听讼的人显得偏袒，判定无辜者有罪、有罪者无辜，则他们不得再听另一件讼案：\Emph{以此}接受他们偏袒的责备，是适宜的。\par

23. 宗徒们进一步规定：高傲自大、被夸耀的傲慢所抬举的人，不得被接纳进入任职：因为\Emph{经上记载}：\Quote{人在人前尊高的，在天主前是可憎的；}又论到他们说：\Quote{我必报复那些自夸的人。}\par

24. 宗徒们进一步规定：在村庄中应有监督管辖长者，并应被承认为众人之首，众人均要在他的手下被追究：因为撒慕尔也曾这样巡行各地并施行治理。\par

25. 宗徒们进一步规定：那些将来相信基督的君王，应准予他们上前，站在祭台前，与教会的指导者们一同站立：因为达味和与他相似的人也曾上前站在祭台前。\par

26. 宗徒们进一步规定：任何人不得凭司祭的权柄做任何不符合正义和公平的事，而应按照正义，免于偏袒之责，\Emph{方可行一切事。}\par

27. 宗徒们进一步规定：供品的饼应在烤制当日置于祭台上，不得数日之后——此事是不允许的。\par

所有这些都是宗徒们规定的，并非为他们自己，而是为以后的人——因为他们担心将来会有披着羊皮的豺狼：因为他们自己身上的圣神、护慰者已足够：正如祂借他们的手制定了这些法律，\Emph{照样}祂会合法地引导他们。因为他们已从我们的主领受了权能和权威，无需别人为他们制定法律。保禄和弟茂德在叙利亚和基里基雅各地巡游时，也将宗徒和长老的同一诫命和法律，交付给那些在宗徒权下的人，为他们宣讲和传播福音的各地方教会。\par

门徒们在他们制定了这些规条与律法之后，并未停止宣讲福音，也未停止主藉他们之手所行的奇妙奇能。因为每天都有许多百姓聚集在他们周围，相信基督；还有人从别的城市来到他们那里，听见他们的话并领受了。尼苛德摩和加玛里耳，犹太会堂的首领，也常暗中来到宗徒们跟前，赞同他们的教导。此外，阿纳尼亚的儿子们——犹达、肋未、佩里、若瑟和犹斯托，以及司祭盖法和亚历山大，也常在夜间来到宗徒们那里，宣认基督是天主子；但他们惧怕本国人民，所以不敢向门徒们表明自己的心意。\par

宗徒们亲切地接待了他们，对他们说：你们不要因人的羞耻和畏惧而丧失在天主面前的救恩，也不要使基督的血归到你们身上，正如你们的祖先曾把它揽在自己身上一样；因为在天主面前，这是不可接受的：你们既然\Emph{暗中}与他崇拜者在一起，却又去与杀害他可敬爱圣子的凶手来往。你们怎能期望你们的信仰被那些真诚的人接纳，而你们却与那些虚伪的人为伍？你们作为相信基督的人，应当公开宣认我们所传的这个信仰。\par

司祭的儿子们听见门徒们所说的这些话后，全都一同在众宗徒面前喊道：我们宣认并相信被钉十字架的基督，我们宣认他从永远就是天主子；我们弃绝那些胆敢钉死他的人。因为百姓中的司祭暗中宣认基督，但他们为了所贪爱的百姓中的首长地位，不愿公开宣认；他们忘记了经上所记载的：\Quote{祂是知识的掌管者，诡诈的计谋在祂面前毫无用处。}\par

他们的父亲从儿子们口中听见这些事，就对他们极其敌视；这并非因为他们信了基督，而是因为他们在本国百姓的儿子们面前公开宣告并说出了他们父亲的心意。\par

但那相信的人都紧贴着门徒们，不离开他们，因为他们看见门徒们所教导众人的一切，都当着众人的面亲自付诸实行；当门徒们遭受磨难和迫害时，他们乐愿与他们一同受苦，并因宣认对基督的信仰而甘受鞭打和监禁；他们一生所有的日子都在犹太人和撒玛黎雅人面前宣讲基督。\par

宗徒们去世之后，教会中有导师和治理者；宗徒们托付给他们、他们从宗徒们领受的一切，他们一生都继续教导众人。他们自己也在临终时，把从宗徒们所领受的一切托付并交付给他们的门徒；还有雅各伯从耶路撒冷写的，西满从罗马城写的，若望从厄弗所写的，马尔谷从大亚历山大里亚写的，安德肋从夫黎基雅写的，路加从马其顿写的，以及犹达多默从印度写的：使宗徒的书信可以在各处的教会中被接纳和诵读，正如路加所写的他们的《宗徒大事录》所记载的事迹被诵读一样；这样，就可以由此认识宗徒们和先知们，以及旧约和新约；从而可以\Emph{看出}，一个真理在他们每一个人里面被宣告出来：一个圣神在他们每一个人里面说话，来自同一位天主，他们都曾崇拜并宣讲过同一位天主。各个\Emph{不同}的国家都领受了他们的教导。因此，我们的主藉宗徒们所说的一切，以及宗徒们交付给门徒们的一切，都在每个国家被相信和接纳，这是藉着我们的主的作为，祂曾对他们说：\BibleQuote{我与你们天天在一起，直到今世的终结。}导师们用先知书与犹太人辩论，并且用他们因基督之名所行的可畏奇能与受迷惑的外邦人争辩。因为所有民族，甚至那些住在遥远国度的人民，都安静地领受了基督的福音；那些成为宣信者的人在迫害中喊道：我们今天所受的迫害将为我们辩护，\Emph{免得我们受惩罚}，因为我们\Emph{自己}先前曾是迫害者。因为他们中有些人被下令处以刀剑之死，而有些人则被剥夺了一切所有后释放。他们遭受的磨难越大，他们的会众就变得越发丰富广大；他们心怀喜乐地接受各种死亡。藉由司铎圣职的授任，也就是宗徒们自己从我们主那里领受的，他们的福音迅速传遍世界的四方。并且通过相互探访，他们彼此服侍。\par

1. 耶路撒冷，以及整个巴勒斯坦地区、撒玛黎雅人居住的地区、培肋舍特人居住的地区、阿拉伯地区、腓尼基地区、凯撒勒雅的人民，都从雅各伯领受了司铎圣职的授任——他是建在熙雍的宗徒教会的治理者和导师。\par

2. 大亚历山大里亚、底比斯、整个内埃及、整个佩鲁西翁地区，以及\Emph{一直延伸}到印度人边界的地方，都从福音作者马尔谷领受了宗徒的司铎圣职授任——他是那里他所建立教会的治理者和导师，并在那里\Emph{履职}。\par

3. 印度以及所有属于它和它周围的地区，直到最远的海洋，都从犹达多默领受了宗徒的司铎圣职授任——他是那里他所建立教会的导师和治理者，并在那里\Emph{履职}。\par

4. \Person{安提约基雅}、\Place{叙利亚}、\Place{基里基雅}、\Place{迦拉达}，直至\Place{本都}，从\Person{西满刻法}接受了宗徒的司铎圣秩，他亲自在那里为教会奠基，并担任司铎、履行职务，直到他因行邪术的西满（此人正以妖术迷惑罗马民众）而离开那里，上到\Place{罗马}。\par

5. \Place{罗马}城、整个\Place{意大利}、\Place{西班牙}、\Place{不列颠}、\Place{高卢}，以及周围所有其余地区，从\Person{西满刻法}接受了宗徒的司铎圣秩；他从\Place{安提约基雅}上去，在那里他建立的教会及其周围地方担任治理者和导师。\par

6. \Place{厄弗所}、\Place{得撒洛尼}、整个\Place{亚细亚}、\Place{科林托}全境、整个\Place{阿哈雅}及其周围地区，从福音者\Person{若望}接受了宗徒的司铎圣秩；他曾依靠在我们主的胸怀，亲自在那里建立教会，并在那里以他所\Emph{担任}的导师职份履行职务。\par

7. \Place{尼凯亚}、\Place{尼科美底亚}、整个\Place{彼提尼雅}地区、\Place{内迦拉达}及其周围区域，从\Person{西满刻法}的兄弟\Person{安德肋}接受了宗徒的司铎圣秩；他亲自在他建立的教会中担任导师和治理者，并担任司铎、履行职务。\par

8. \Place{拜占庭}、整个\Place{色雷斯}地区及其周围地带，直至那条作为与蛮族分界的大河，从宗徒\Person{路加}接受了宗徒的司铎圣秩；他亲自在那里建立教会，并以他所\Emph{担任}的治理者及导师职份履行职务。\par

9. \Place{厄德撒}及其四围所有地区、\Place{祚巴}、\Place{阿拉伯}、整个北方及其周围区域、南方以及\Place{美索不达米亚}边境所有地区，从七十二宗徒之一的宗徒\Person{阿戴}接受了宗徒的司铎圣秩；他亲自在那里收门徒，建立教会，并担任司铎、履行职务，以他所\Emph{担任}的导师职份。\par

10. 整个\Place{波斯}、\Place{亚述人}、\Place{亚美尼亚人}、\Place{玛待人}、\Place{巴比伦}周围地区、\Place{胡兹人}和\Place{格莱人}，直至\Place{印度人}的边界，直至\Place{哥格和玛哥格}之地，此外还有四围所有国家，从\Person{阿盖}——\Person{阿戴}宗徒的门徒、一名丝绸工匠——接受了宗徒的司铎圣秩。\par

宗徒其余的同伴又前往遥远的蛮族地区；他们随处收门徒，继续前行；在那里以他们的宣讲履行职务；他们的离世也发生在那里；他们的门徒在他们之后继续\Emph{这项工作}直到今日，他们在宣讲中既未作任何更改，也未有所添加。\par

此外，福音者\Person{路加}如此勤勉，以致他写下了《\WorkTitle{宗徒大事录}》的业绩，以及他们司铎职份的规条和法律，以及他们各人所去之处。我说，藉着他的勤勉，路加写下了这些事，甚至更多；他将这些交托给他的门徒\Person{普黎斯库}和\Person{阿桂禄}；他们一直陪伴他直到他去世之日，正如\Person{吕斯特辣的厄辣斯托}和\Person{默纳乌}——这些宗徒的首批门徒——以及\Person{弟茂德}陪伴\Person{保禄}，直到他因对抗演说家\Person{特图禄}而被押解至\Place{罗马}城。\par

\Person{尼禄凯撒}在\Place{罗马}城用刀剑处死了\Person{西满刻法}。\par

\SourceRecord{Translated by B.P. Pratten.；Ante-Nicene Fathers；Vol. 8.；Edited by Alexander Roberts, James Donaldson, and A. Cleveland Coxe.；Buffalo, NY: Christian Literature Publishing Co.,；1886.；<http://www.newadvent.org/fathers/0854.htm>.}
